\noindent
This thesis presents the design, construction, and experimental characterization of a quartz-enhanced photoacoustic spectroscopy (QEPAS) sensor for the detection of trace nitrous oxide (\nto). The sensor uses a pulsed, tunable quantum cascade laser (QCL) as the excitation source, operated in amplitude-modulation mode at the resonance frequency of a quartz tuning fork housed in a commercial acoustic detection module, with the photoacoustic signal recovered by a lock-in amplifier. The laser is tuned to the \nto\ $\nu_1$ band near \SI{7.7}{\micro\meter}, a line chosen to lie within the tuning range of the available source. The work covers the integration and alignment of the optical, acoustic, and electronic subsystems, an external-trigger amplitude-modulation scheme that includes a solution to the loss of lock at short pulse widths, and a wavelength-specific optical power budget for the beam-delivery path. The tuning fork, the lock-in detection chain, and the laser are characterized individually, and \nto\ is detected and calibrated against concentration in order to quantify the sensor performance in terms of sensitivity, normalized noise-equivalent absorption, minimum detectable concentration, and stability from an Allan-deviation analysis. \todo{Insert the final headline figures (NNEA, minimum detectable concentration, optimal integration time) once the measurements are complete.} The resulting instrument provides a validated platform intended for the future detection of trace volatile organic compounds.
\\

\noindent
\textbf{\textit{Keywords}}: quartz-enhanced photoacoustic spectroscopy (QEPAS), nitrous oxide, pulsed quantum cascade laser, amplitude modulation, trace gas detection
